\subsection{Summary of Findings}

This paper addressed day-ahead electricity demand forecasting under asymmetric
operational costs.  Our main findings are:

\begin{enumerate}
    \item \textbf{Asymmetric costs are pervasive and substantial.}  In
          electricity markets the costs of under-predicting demand substantially
          exceed those of over-prediction.  Empirically plausible cost ratios
          $c_u/c_o\in[5,20]$ have a large impact on the optimal forecasting
          strategy.

    \item \textbf{The newsvendor framework provides a principled optimality
          condition.}  The cost-minimising procurement level corresponds to a
          specific quantile $\tau^*=c_u/(c_u+c_o)$ of the conditional demand
          distribution.  For a cost ratio of~5 this yields $\tau^*\approx 0.833$,
          i.e.\ the 83.3rd percentile of demand --- a systematic upward shift
          that is normatively \emph{desirable} under asymmetric costs.

    \item \textbf{Asymmetric loss training effectively targets the cost-minimising
          quantile.}  A feedforward neural network trained with a cost-sensitive asymmetric loss
          function (specifically, a non-linear variant of the pinball loss) at
          the appropriate quantile achieves lower expected operational cost than
          an L1-trained neural network baseline.  The cost
          reduction is consistent across a wide range of synthetic cost function
          parameters.

    \item \textbf{Symmetric loss is a misleading evaluation metric under asymmetric costs.}
          The L1-trained neural network achieves lower MAE than the asymmetric-loss model
          yet incurs higher expected operational cost.  Practitioners should adopt
          cost-sensitive metrics such as asymmetric loss or expected operational
          cost for model selection and evaluation.

    \item \textbf{Post-hoc calibration helps but does not substitute for
          cost-aware training.}  A global affine correction fitted to the
          operational cost recovers much of the gap on asymmetric metrics, but
          cannot match a model trained end-to-end, because it cannot adapt the
          upward bias to the input-conditional structure of forecast error.

    \item \textbf{The approach is practical and improves on the operational
          forecast.}  Our model is a small, open-source feedforward network that
          trains in minutes on commodity hardware.  Even though predictive
          accuracy was not an objective of this work, it improves on the Noga
          operational forecast not only on cost-sensitive metrics but even under
          symmetric error metrics such as MSE --- and we believe substantial
          further gains are readily achievable.  Together with a demonstrated
          economic benefit, this makes the method realistic to adopt in practice.
\end{enumerate}

\subsection{Implications for Market Participants and Policymakers}

\textbf{For market participants} (electricity retailers, large consumers,
aggregators), the key takeaway is that standard forecasting tools optimised for
MSE are likely to systematically under-procure electricity, exposing them to
avoidable real-time market risk.  Incorporating asymmetric cost modelling ---
through direct asymmetric-loss training or post-hoc quantile adjustment --- can
yield meaningful reductions in procurement costs with negligible additional
computational overhead.

\textbf{For grid operators} (transmission system operators, independent system
operators), the implication is that incentivising cost-sensitive forecasting
among market participants could improve overall system reliability.  If all
participants systematically under-forecast demand, the system operator faces
persistent shortfalls in committed reserves.  Regulatory frameworks that reward
forecast accuracy measured by cost-sensitive metrics, rather than MSE, would
align private incentives with system reliability objectives.

\textbf{For policymakers}, our results support the inclusion of cost-sensitive
forecasting requirements in market design rules.  For example, mandatory reporting of asymmetric (pinball) loss at relevant
quantiles, alongside MSE, would provide a richer picture of forecast quality
that better reflects operational costs.

\subsection{Limitations}

\begin{itemize}
    \item \textbf{Synthetic cost functions.}  Since we do not have access to
          actual cost data from the Israeli electricity market, all cost
          evaluations rely on synthetic functions that satisfy qualitative
          constraints.  The results may not precisely reflect actual market costs.

    \item \textbf{Time-invariant cost asymmetry.}  Although our model operates at
          fine (5-minute) granularity, we treat the cost asymmetry as constant
          across the day.  In practice, markets clear at sub-hourly intervals and
          the asymmetry of costs may vary significantly across hours of the day.

    \item \textbf{Point forecasts only.}  Our neural network produces a single scalar
          forecast at the target quantile.  A richer approach would estimate the
          full predictive distribution, enabling uncertainty quantification and
          adaptive quantile selection.

    \item \textbf{Fixed cost ratio.}  We assume that $c_u/c_o$ is constant over
          time.  In practice the relative costs may vary with fuel prices, demand
          levels, and seasonal reserve availability.
\end{itemize}

\subsection{Directions for Future Research}

\begin{enumerate}
    \item \textbf{Distributional forecasting.}  Training models to output full
          predictive distributions (via normalising flows, conformal prediction,
          or Bayesian deep learning) would enable the optimal quantile to be
          selected adaptively as a function of current cost parameters.  We regard
          this, rather than training against a fixed cost function, as the proper
          long-term formulation of the problem, for the reasons set out in the
          concluding remarks below.

    \item \textbf{Nonlinear and non-smooth cost functions.}  Step-function or
          tiered cost structures require optimisation approaches beyond the
          pinball loss, such as distributionally robust optimisation or stochastic
          programming.

    \item \textbf{Online and adaptive learning.}  Adaptive algorithms that update
          the forecasting model in real time as market conditions change would be
          valuable for practical deployment.

    \item \textbf{Richer temporal modelling and time-varying cost asymmetry.}
          Our model already operates at fine (5-minute) granularity but treats
          each target slot largely independently and assumes a constant cost
          ratio.  Explicitly modelling intra-day and cross-day demand dynamics
          --- for example with Transformer
          architectures~\parencite{vaswani2017attention} or temporal
          convolutional networks --- together with a cost asymmetry that varies
          by hour of the day, could better exploit the structured hour-of-week
          error pattern documented in Section~\ref{sec:forecast-error-dist}.

    \item \textbf{Multi-product forecasting.}  In practice, electricity retailers
          must manage procurement across multiple products (day-ahead energy,
          reserves, ancillary services) with correlated costs and demands.  A
          joint cost-sensitive forecasting framework would be a substantial
          extension.

    \item \textbf{Real cost data.}  Collaboration with market participants or
          operators to obtain actual cost data would enable precise calibration
          and remove the need for synthetic cost functions.
\end{enumerate}

\subsection{Concluding Remarks}

Electricity demand forecasting sits at the intersection of time series analysis,
optimisation under uncertainty, and energy economics.  As markets become more
complex --- driven by variable renewable energy, demand response, and distributed
storage --- the costs of forecast errors are likely to increase, making
cost-sensitive forecasting ever more important.

This work demonstrates that a conceptually simple modification --- replacing the
MSE training objective with a cost-sensitive asymmetric loss function
(specifically, a non-linear variant of the pinball loss at the cost-minimising
quantile) --- leads to meaningful improvements in operational performance without
sacrificing computational tractability.

We stress, however, that training against a \emph{predefined, fixed} cost
function is not the solution we ultimately advocate.  It couples the learned
model to one particular cost structure, so that a change in the cost ratio ---
driven by fuel prices, reserve availability, or the time of day --- in principle
requires retraining, and it collapses the rich uncertainty in future demand into
a single quantile chosen in advance.  We view it as one useful step in the right
direction rather than the destination.  The desired formulation, in our view, is
to \emph{decouple estimation from decision-making}: first predict the full
\emph{conditional distribution} of demand given the available features, and then,
at decision time, optimise the procurement quantity against whatever cost
function currently prevails.  This separation would let a single trained model
serve any cost structure, adapt the optimal bias to the input-conditional
uncertainty of each slot, and expose the full distribution for downstream
risk-aware decisions.  Realising it is the main direction of our future work.

We hope these findings contribute to a broader shift
towards cost-aware forecasting practices in the electricity industry and stimulate
further research at the intersection of machine learning and energy systems
optimisation.
